\chapter{Conclusiones}
\label{cap:conclusiones}

\section{Resumen del trabajo}

Este trabajo ha consistido en implementar desde cero un entorno completo para el análisis de generadores pseudoaleatorios basados en LFSR
y usarlo para comparar cuatro construcciones clásicas, Geffe, Shrinking, Self-Shrinking y Alternating Step Generator.

El entorno incluye tres componentes principales. El primero es la implementación de los generadores, un LFSR genérico parametrizable por grado y polinomio,
y los cuatro generadores construidos sobre él. El segundo es el algoritmo de Berlekamp-Massey,
que permite medir la complejidad lineal de cualquier secuencia binaria y calcular su perfil de evolución.
El tercero son ocho tests del estándar NIST SP 800-22, cada uno evaluando un aspecto distinto de la aleatoriedad
(balance, distribución de rachas, rango matricial, compresibilidad, distribución de patrones, espectro de frecuencias y complejidad lineal por bloques).

Son unas 1500 líneas de Python escritas sin librerías criptográficas. Las únicas dependencias externas son scipy (funciones de distribución estadística),
numpy (operaciones con arrays) y matplotlib (gráficas). Hay 27 tests unitarios que verifican cada componente contra valores de referencia conocidos.

El análisis experimental se ha realizado con secuencias de 500.000 bits para los tests NIST y 5.000 bits para la complejidad lineal,
con un nivel de significación $\alpha = 0{,}01$. Los resultados principales son los siguientes,
el generador de ASG pasa los 8 tests con complejidad lineal cercana a $N/2$,
Shrinking y Self-Shrinking tienen complejidad lineal similar al ASG pero fallan el test espectral,
el generador de Geffe pasa 7 tests pero es vulnerable a ataques de correlación, y el LFSR puro falla los 3 tests que detectan estructura lineal.

\section{Conclusiones principales}

\subsection{La complejidad lineal es necesaria pero no suficiente}

La complejidad lineal separa claramente a los generadores en dos grupos. Por un lado el LFSR puro y Geffe,
con valores bajos respecto a la longitud de la secuencia. Por otro Shrinking, Self-Shrinking y ASG, con valores cercanos a $N/2$.
Sin embargo, dentro del segundo grupo la complejidad lineal no distingue la calidad real.

Geffe tiene complejidad lineal 623, mucho mayor que la del LFSR puro (19) pero baja en términos absolutos.
Aun así, su debilidad principal no es esa sino la correlación $0{,}75$ con dos de sus LFSR internos,
que permite un ataque con coste proporcional a la suma de los espacios de estados, no al producto.
Un atacante rompe Geffe mucho antes de lo que Berlekamp-Massey requeriría.

Los otros tres generadores tienen complejidades lineales casi idénticas ($\sim 2500$) pero comportamiento distinto en el test espectral,
y eso la complejidad lineal no lo captura. Tener LC alta es condición necesaria (un generador con LC baja cae frente a Berlekamp-Massey)
pero no basta para garantizar buenas propiedades estadísticas.

\subsection{Los tests estadísticos detectan propiedades, no seguridad}

Los tests NIST detectan desviaciones estadísticas en la secuencia de salida, no evalúan la seguridad del generador como construcción criptográfica.
Geffe pasa 7 de 8 tests pese a ser trivialmente atacable. El LFSR puro pasa 5 de 8 pese a ser completamente predecible con 38 bits observados.

Esto no significa que los tests sean inútiles. Si la salida tiene patrones detectables, el generador no puede ser seguro.
Pero pasar todos los tests no garantiza seguridad, solo que no hay defectos estadísticos evidentes.
Los tests filtran generadores con defectos obvios, no sustituyen al análisis de la estructura interna.

\subsection{El mecanismo de combinación determina la calidad}

De los cuatro generadores, el ASG es el que mejor se comporta tanto en complejidad lineal como en los tests estadísticos.
Pasa los 8 tests NIST con los mismos parámetros con los que los otros tres fallan al menos uno.
La diferencia está en su mecanismo de combinación, reloj irregular más XOR de dos fuentes independientes.

La comparación más clara es entre Geffe y ASG, ya que ambos usan exactamente los mismos tres LFSR.
La única diferencia es cómo se conectan. Geffe usa un multiplexor (selección directa con correlación) y ASG usa reloj alterno más XOR.
El resultado, Geffe tiene LC 623 y falla espectralmente, ASG tiene LC 2499 y pasa todo.
Mismo hardware, resultado completamente distinto.

La conclusión es clara, en generadores basados en LFSR la clave no es el tamaño de los registros sino cómo se combinan.
Un buen mecanismo convierte tres LFSR pequeños en un generador robusto, un mecanismo malo (el multiplexor de Geffe) deja el sistema vulnerable
da igual cuánto se agranden los registros.

\subsection{La selección irregular es buena pero no perfecta}

Shrinking y Self-Shrinking logran complejidades lineales excelentes y pasan 7 de 8 tests.
Su mecanismo de selección irregular destruye la estructura lineal de la secuencia base.
Sin embargo, ambos fallan el test espectral, lo que indica que la periodicidad del LFSR base no se elimina completamente con el muestreo.

Que el ASG (con un mecanismo distinto) sí pase el espectral sugiere que para eliminar las componentes frecuenciales
no basta con muestrear irregularmente una secuencia periódica, hace falta combinar dos fuentes independientes de forma que sus periodicidades
se cancelen mutuamente.

En la práctica esta debilidad espectral probablemente desaparece con LFSR de grado suficiente,
porque los armónicos se diluyen en secuencias largas. Pero con los grados pequeños usados en este experimento, la diferencia es medible.

\section{Cumplimiento de objetivos}

Se han cumplido todos los objetivos de la introducción.

\begin{enumerate}
  \item Se ha implementado un LFSR genérico y el algoritmo de Berlekamp-Massey desde cero en Python,
        con tests unitarios que verifican su corrección.
  \item Se han implementado los cuatro generadores (Geffe, Shrinking, Self-Shrinking y ASG),
        todos funcionales y verificados contra resultados conocidos.
  \item Se ha medido la complejidad lineal de cada generador,
        verificando que los valores experimentales son coherentes con las cotas teóricas.
  \item Se han pasado 8 tests NIST a las secuencias de cada generador con $\alpha = 0{,}01$,
        y se han comparado los resultados.
  \item Se han extraído conclusiones sobre la seguridad relativa de los generadores,
        justificadas tanto por los resultados experimentales como por la teoría.
\end{enumerate}

Además de estos objetivos explícitos, el trabajo ha producido un software completo y documentado que podría usarse como base para análisis más amplios
(con más tests, más generadores o parámetros más grandes).

\section{Trabajo futuro}

\subsection{Secuencias más largas y más estados iniciales}

El análisis usa un estado inicial fijo por generador. Una extensión natural es ejecutar los tests con múltiples estados iniciales seleccionados al azar
y reportar la distribución de p-valores en vez de un único valor.
Así se podría ver si los resultados cercanos al umbral (como el block frequency del Self-Shrinking con 0,012) son estables o dependen del estado concreto.

Otra dirección interesante es fijar la longitud de secuencia y variar los grados de los LFSR para trazar la frontera de detección,
a partir de qué grado mínimo el Shrinking empieza a pasar el test espectral.
El resultado sería una recomendación práctica de tamaño mínimo para cada generador.

\subsection{Implementación de ataques}

El ataque de correlación contra Geffe se ha descrito teóricamente pero no se ha implementado.
Una extensión interesante sería programarlo y verificar experimentalmente que el coste real coincide con el predicho ($2^{n_1} + 2^{n_2} + 2^{n_3}$).
Dado un fragmento de keystream observado, probar los $2^5 = 32$ estados posibles de $R_2$, generar la secuencia correspondiente y contar coincidencias.
El estado correcto debería dar coincidencia $\sim 75\%$ y uno incorrecto $\sim 50\%$.
Verificar esto en la práctica y medir cuántos bits necesita la distinción para ser fiable sería una demostración concreta de la vulnerabilidad.

También se podría implementar fuerza bruta sobre $R_1$ del Shrinking y medir cuántos bits son necesarios para identificar correctamente la secuencia de selección.

\subsection{Tests adicionales y herramientas externas}

Implementar los 7 tests NIST restantes daría una imagen más completa del comportamiento de cada generador.
Random excursions en particular podría revelar asimetrías que los implementados no detectan.

Para validar la implementación se podría ejecutar la batería NIST oficial (escrita en C por el propio NIST) sobre las mismas secuencias generadas por mi código.
Esto serviría como comprobación cruzada, y si los p-valores coinciden con la implementación oficial,
se tendría confianza adicional en la corrección.

\subsection{Generadores más recientes}

Los cuatro generadores estudiados son construcciones clásicas de los años 70-90.
Sería interesante repetir el análisis con generadores más recientes como Grain o Trivium (del proyecto eSTREAM),
que también están basados en registros de desplazamiento pero incorporan componentes no lineales más sofisticados.
Comparar estos con los clásicos permitiría ver cuánto ha avanzado el diseño de cifrados de flujo en las últimas décadas.

Grain combina un LFSR con un NLFSR (registro con retroalimentación no lineal) del mismo tamaño y aplica una función booleana de filtrado sobre ambos.
Implementar Grain con las mismas herramientas de este trabajo y compararlo con el ASG mostraría si el componente no lineal explícito da mejores resultados
que la no linealidad implícita del reloj irregular.

\section{Valoración personal}

Este trabajo me ha permitido entender la criptografía de flujo a un nivel de detalle que no habría conseguido solo leyendo sobre ella.
Implementar Berlekamp-Massey desde cero, por ejemplo, obliga a entender por qué funciona y cómo cada paso del algoritmo contribuye al resultado.
Lo mismo pasa con los tests NIST, ya que programar cada uno te fuerza a pensar en qué hipótesis estadística se está contrastando
y qué distribución teórica se asume.

La parte que me resultó más sorprendente fue que el ASG (con los mismos LFSR que el de Geffe) diera resultados tan superiores en los tests estadísticos.
Antes de hacer los experimentos esperaba que Shrinking y Self-Shrinking fueran los mejores por ser los más recientes.
Que el de ASG (de 1987, anterior a los otros dos) fuera el único en pasar todos los tests fue contraintuitivo hasta entender la razón,
el XOR con reloj alterno destruye la periodicidad de una forma que el muestreo irregular no consigue.

También me pareció instructivo ver que el LFSR puro pasa 5 de 8 tests. Si solo se evaluara con tests estadísticos básicos parecería un generador aceptable.
Esto refuerza algo que leí varias veces pero no terminaba de interiorizar, las propiedades estadísticas y la seguridad criptográfica son cosas distintas.
Para evaluar un generador hay que atacarlo, no solo testearlo.

Me llevo el haber aprendido a usar un asistente de inteligencia artificial como apoyo para revisar el código y la redacción.
Esto vino a raiz de que en mi plano profesional nos animaban a usar estas herramientas de manera constructiva, ya que se observaba una gran mejora en la calidad del trabajo cuando se usaban correctamente,
ademas, de un claro beneficio en el tiempo de trabajo, ya que uno empieza a usarlas como un colaborador que revisa tu trabajo, y lo emplea para contrastar, encontrar erratas o cuestionar una explicación poco clara, revisando siempre lo que devuelve.
Es una herramienta que en el entorno laboral se pide cada vez más, y saber usarla de forma crítica (sin delegar en ella el criterio) me parece una competencia tan valiosa como las
puramente técnicas del trabajo.

Por último, la decisión de implementar todo desde cero ha sido la que más tiempo ha costado pero también la que más me ha enseñado.
Programar el test espectral, por ejemplo, me obligó a entender la DFT a un nivel que de otra forma no habría necesitado.
Y cuando el Self-Shrinking dio un p-valor de 0,012 en block frequency, pude ir directamente al código del test a comprobar que no era un bug
(no lo era, simplemente esa secuencia concreta tiene un bloque con desequilibrio local). Esa capacidad de ir al código y entender exactamente qué pasa
es lo que justifica el esfuerzo de no usar librerías externas.
